\section{Appendix}

%%%%GPT Writing Assistance 
\begin{lemma}\label{lemma:conv_ineq}
    Let $m \in \mathbb{Z}^2$ and $\alpha \in (1, 2)$. Then we have the bound
    \begin{equation*}
        S(m) \coloneqq \sum_{n \in \mathbb{Z}^2} \langle n \rangle^{-\alpha} \langle n + m \rangle^{-\alpha} \lesssim_\alpha \langle m \rangle^{2-2\alpha}.
    \end{equation*}
    If $\alpha = 2$, then we have the bound
    \begin{align*}
        S_2(m) \coloneqq \sum_{n \in \mathbb{Z}^2} \langle n \rangle^{-2} \langle n + m \rangle^{-2} \lesssim \jb{m}^{-2} (1+\ln(\jb{m}^2)).
    \end{align*}
\end{lemma}

\begin{proof}
    We cover the domain of summation into three regions: $R_1 \coloneqq \{n \in \mathbb{Z}^2 : |n| \le \frac{|m|}{2}\}$, $R_2 \coloneqq \{n \in \mathbb{Z}^2 : |n+m| \le \frac{|m|}{2}\}$, and $R_3 \coloneqq \{n \in \mathbb{Z}^2 : |n| > \frac{|m|}{2} \text{ and } |n + m| > \frac{|m|}{2}\}$. Since $\mathbb{Z}^2 \subset R_1 \cup R_2 \cup R_3$, establishing the bound on each of these three subregions is sufficient to prove the statement for the entire sum.

    The bound for $R_1$ follows by applying the reverse triangle inequality. For any $n \in R_1$, we have $|n| \le \frac{|m|}{2}$, which implies:
    \begin{equation*}
        |n+m| \ge |m| - |n| \ge |m| - \frac{|m|}{2} = \frac{|m|}{2}.
    \end{equation*}
    Because $\alpha > 0$, this implies $\langle n + m \rangle^{-\alpha} \lesssim \langle m \rangle^{-\alpha}$. Approximating the remaining sum by a polar integral and utilizing the fact that $2 - \alpha > 0$ for $\alpha \in (1, 2)$, we obtain
    \begin{align}\label{S_R1_bound}
        \begin{split}
        S_{R_1}(m) &= \sum_{n \in R_1} \langle n \rangle^{-\alpha}\langle n+m \rangle^{-\alpha} 
        \lesssim \langle m \rangle^{-\alpha} \sum_{|n| \le \frac{|m|}{2}} \langle n \rangle^{-\alpha} \\
        &\lesssim \langle m \rangle^{-\alpha} \left(1+\int_0^{\frac{|m|}{2}} (1 + r^2)^{-\frac{\alpha}{2}} r \, dr \right)
        = \langle m \rangle^{-\alpha} \left(1 + \frac{1}{2 - \alpha} \left[ (1+r^2)^{1-\frac{\alpha}{2}} \right]_{0}^{\frac{|m|}{2}} \right) \\
&\lesssim_\alpha \langle m \rangle^{-\alpha}
    \left(1+\langle m \rangle^{2-\alpha}\right)
\lesssim_\alpha \langle m \rangle^{2-2\alpha}.        
\end{split}
    \end{align}
    The bound in region $R_2$ follows by shifting the summation index to $n' = n+m$. 

    To show the bound in the final region $R_3$, we apply Young's product inequality
    \begin{equation*}
        S_{R_3}(m) \le \frac{1}{2} \sum_{n \in R_3}\langle n \rangle^{-2\alpha} + \frac{1}{2} \sum_{n \in R_3}\langle n+m \rangle^{-2\alpha}.
    \end{equation*}
    By symmetry, it suffices to bound the first term. Comparing the sum to a polar integral and noting that $1 - \alpha < 0$, we obtain
    \begin{align*}
        \sum_{n \in R_3} \langle n \rangle^{-2\alpha} 
        \le \sum_{|n| > \frac{|m|}{2}} \langle n \rangle^{-2\alpha} 
        \lesssim \int_{\frac{|m|}{2}}^\infty (1+r^2)^{-\alpha} r \, dr 
        = \frac{1}{2 - 2\alpha} \left[ (1+r^2)^{1-\alpha} \right]_{\frac{|m|}{2}}^\infty 
        \lesssim \langle m \rangle^{2-2\alpha}.
    \end{align*}
    The proof for the bound of $S_2$ follows analogously, but the estimate in \eqref{S_R1_bound} changes to 
    \begin{align*}
        S_{2_{R_1}}(m) &= \sum_{n \in R_1} \langle n \rangle^{-2}\langle n+m \rangle^{-2} 
        \\
        &\lesssim \langle m \rangle^{-2} \left(1+ \int_0^{\frac{|m|}{2}} (1 + r^2)^{-1} r \, dr \right)
        \\
%&= \langle m \rangle^{-2}\left(1+\frac12\ln\left(1+\frac{|m|^2}{4}\right)\right) \\
&\leq \langle m \rangle^{-2}
    \left(1+\ln(1+|m|^2)\right).
    \end{align*}
    This yields the desired estimate.
\end{proof}
%%%%

%%%%GPT alone:
% Next, we demonstrate that $\H^{-\sigma}(\T^2)$ is the natural support space of any invariant measure of the \eqref{eq:SDMFW}.
% \begin{proposition}[Sobolev regularity of the stationary Gaussian law]
% \label{prop:stationary_sobolev_regularity}
% Let $\eps>0$ and $\Lambda=1-\Delta$ on $\mathbb{T}^2$.
% Suppose that $L$ is a strictly positive self-adjoint operator on
% $L^2(\mathbb{T}^2)$ with form domain $H^1(\mathbb{T}^2)$ and that
% \[
%     a\Lambda \leq L \leq b\Lambda
% \]
% in the sense of quadratic forms, for some $a,b>0$.
% Let $\nu$ be the centered Gaussian law on
% $\mathcal{D}'(\mathbb{T}^2)\times\mathcal{D}'(\mathbb{T}^2)$
% with covariance
% \[
%     \begin{pmatrix}
%         L^{-1} & 0\\
%         0 & \eps^{-2} I
%     \end{pmatrix},
% \]
% where covariance is understood through pairings with smooth test
% functions.

% Then
% \[
%     \nu\left(
%         \bigcap_{\delta>0}
%         \left[
%             H^{-\delta}(\mathbb{T}^2)
%             \times H^{-1-\delta}(\mathbb{T}^2)
%         \right]
%     \right)=1.
% \]
% Moreover, both endpoint carry no mass:
% \[
%     \nu\bigl(\{(u,v):u\in L^2(\mathbb{T}^2)\}\bigr)=0,
%     \qquad
%     \nu\bigl(\{(u,v):v\in H^{-1}(\mathbb{T}^2)\}\bigr)=0.
% \]
% \end{proposition}

% \begin{proof}
% Let $(e_j)_{j\geq1}$ be a real orthonormal eigenbasis of $\Lambda$ in $L^2(\mathbb{T}^2)$, with $\Lambda e_j=\lambda_j e_j$. On $\mathbb{T}^2$,
% \[
%     \sum_{j\geq1}\lambda_j^{-1-\delta}<\infty
%     \quad\text{for every }\delta>0,
%     \qquad
%     \sum_{j\geq1}\lambda_j^{-1}=\infty.
% \]
% The form inequalities imply
% \[
%     b^{-1}\Lambda^{-1}
%     \leq L^{-1}
%     \leq a^{-1}\Lambda^{-1}.
% \]
% Thus, by Tonelli's theorem, for $(u,v)$ with law $\nu$,
% \begin{align*}
%     \mathbb{E}_{\nu}\|u\|_{H^{-\delta}}^2
%     &=
%     \sum_{j\geq1}\lambda_j^{-\delta}
%         \langle L^{-1}e_j,e_j\rangle_{L^2}
%     \leq
%     a^{-1}\sum_{j\geq1}\lambda_j^{-1-\delta}
%     <\infty,\\
%     \mathbb{E}_{\nu}\|v\|_{H^{-1-\delta}}^2
%     &=
%     \eps^{-2}\sum_{j\geq1}\lambda_j^{-1-\delta}
%     <\infty.
% \end{align*}
% Consequently, $\nu(H^{-\delta}\times H^{-1-\delta})=1$ for every $\delta>0$. Taking the countable intersection over $\delta=1/k$, $k\geq1$, proves the first assertion.

% To prove that the endpoints have zero-mass, consider either
% $Y=u$ or $Y=\Lambda^{-1/2}v$.
% The covariance matrix $K_N$ of the real Gaussian vector
% \[
%     \bigl(\langle Y,e_1\rangle,\ldots,
%           \langle Y,e_N\rangle\bigr)
% \]
% satisfies
% \[
%     K_N\geq
%     c\,\operatorname{diag}(\lambda_1^{-1},\ldots,\lambda_N^{-1}),
% \]
% where $c=b^{-1}$ in the first case and $c=\eps^{-2}$
% in the second case. Gaussian integration and monotonicity of
% the determinant on positive definite matrices give
% \begin{align*}
%     \mathbb{E}_{\nu}
%     \exp\left(-\sum_{j=1}^N|\langle Y,e_j\rangle|^2\right)
%     &=
%     \det(I_N+2K_N)^{-1/2}\\
%     &\leq
%     \prod_{j=1}^N(1+2c\lambda_j^{-1})^{-1/2}
%     \longrightarrow 0.
% \end{align*}
% The convergence follows from
% $\sum_j\lambda_j^{-1}=\infty$.
% By bounded convergence, with the convention $e^{-\infty}=0$,
% \[
%     \mathbb{E}_{\nu}
%     \exp\left(-\sum_{j\geq1}|\langle Y,e_j\rangle|^2\right)=0.
% \]
% Hence the series diverges almost surely. Applying this first
% to $Y=u$ and then to $Y=\Lambda^{-1/2}v$ yields the two
% endpoint assertions.
% \end{proof}

